\weeknum{3}
\topic[Inversion]{Exploring \& Understanding Inversion}

\workshopstart

\vspace{-0.75em}
\noindent\textbf{Duration:} 3 hours \hfill \textbf{Setting:} Lecture room with laptops (sphere inversion app)

\section*{Preparation}

\begin{description}
  \item[Pre-reading: Course Notes Ch. 13]
  \item[Lecture videos:]
  \item[Notes:]
\end{description}

% ----------------------------------------------------------
\section*{Purpose}
This workshop introduces the structure of inverse problems in potential-field geophysics, contrasting forward prediction with inverse estimation. Students explore the role of misfit, data weighting, and regularization, and examine how non-linearity and non-uniqueness arise in gravity modelling. Through a sequence of activities, they classify linear vs.\ non-linear models, reason qualitatively about non-uniqueness, and engage in hands-on inversion of a buried sphere using grid search and Gauss--Newton methods.

% ----------------------------------------------------------
\section*{Learning Outcomes}
\begin{itemize}
  \item Distinguish forward and inverse problems and articulate the role of misfit, weighting, and regularization.
  \item Classify simple geophysical relations as linear or non-linear and propose appropriate linearizations.
  \item Recognize manifestations of non-uniqueness in gravity and magnetic profiles and propose discriminating data.
  \item Derive or confirm analytical Jacobian expressions for the buried sphere model.
  \item Interpret misfit surfaces, inversion paths, and the effect of noise and damping.
  \item Explain parameter trade-offs and resolution limitations inherent to the buried-sphere problem.
\end{itemize}

% ----------------------------------------------------------
\section*{Session Flow (\timelinelength\ min)}

Mixed groups of 4--5 with geology and physics/engineering backgrounds.  

\timeline[slot]{Warm-up}{10}

What can you learn from a shadow?
\begin{itemize}
  \item What can you determine directly from the shadow?
  \item What can you estimate only if additional assumptions are made?
  \item What information cannot be determined uniquely?
  \item What prior knowledge are you using when deciding what the object might be?
\end{itemize}

How does this relate to inversion?
\begin{itemize}
  \item A shadow does not uniquely determine the object that created it.
  \item Likewise, a gravity anomaly does not uniquely determine the density distribution that produced it.
  \item The challenge of inversion is deciding which of many acceptable models is most reasonable.
\end{itemize}

\timeline[slot]{Mini-lecture --- Forward vs Inverse, Misfit, Weighting, Challenges}{30}
\begin{itemize}
  \item Forward vs.\ inverse models: $m \mapsto g(m)$ vs.\ $d \mapsto m$.
  \item Misfit: residual $r=d-g(m)$; weighted least squares $\Phi_d=\lVert W_d r\rVert_2^2$.
  \item Role of $W_d$ (noise) and $W_m$ (model constraints); damping $\lambda$.
  \item Challenges: non-linearity (Jacobian needed), non-uniqueness (multiple $m$ give similar $g(m)$), ill-conditioning.
\end{itemize}

\timeline[item]{Testing Linearity}{20}

Two conditions must be met for linearity.  Let's expore how to rigorously test for linearity before we work on other functions.

\timeline[item]{Linear or Non-linear?}{30}
Using the equation handout:
\begin{itemize}
  \item Identify model parameters and classify each relation as \emph{linear}, \emph{non-linear}, or \emph{linearizable}.
  \item If non-linear, propose a linearization (log-transform, re-parameterization, Taylor expansion/Jacobian).
  \item Deliverable: brief notes per group.
\end{itemize}

\setstarttime{70} % 10 min transition/break between activities

% \timeline[slot]{Non-uniqueness}{20}
% Curves-only reasoning:
% \begin{itemize}
%   \item Scenario A: gravity profiles --- two distinct geometric models can reproduce similar anomalies.
%   \item Scenario B: magnetic dike profiles --- width/depth/magnetization trade-offs.
%   \item For \emph{each} curve: propose at least two plausible geologic/geometric models.
%   \item List discriminating data (station density, gradients, susceptibility/remanence, multi-component data, mapping).
%   \item Deliverable: 1-paragraph argument explaining which model to pursue first and what data to collect next.
% \end{itemize}

\timeline[item]{Buried Sphere --- Background}{15}

We're going to look at the gravity anomaly from a buried sphere.  The forward problem is non-linear with respect to the model parameters, requiring a non-linear inverse solution.  This activity is broken into two parts: a look at the formulation, and an investigation using a simple program.
\begin{itemize}
  \item App overview: Data \& Model view; Misfit view; controls for acquisition, noise, synthetic model.
  \item Construct a true model \& generate data (freeze a single noise realization).
\end{itemize}

\timeline[slot]{Buried Sphere --- Jacobian Derivatives}{20}
\begin{enumerate}
  \item Students derive/confirm analytical expressions for
    $\partial g/\partial z$,
    $\partial g/\partial R$,
    $\partial g/\partial \Delta\rho$.
  \item Compare with the app's Jacobian; discuss relative sensitivities:
    amplitude $\propto a^3\Delta\rho$, width controlled by depth $z$.
\end{enumerate}

\timeline[slot]{Buried Sphere --- Investigation}{45}

With the use of \verb+buried_sphere_gui.py+, we will examine a number of aspects of inversion, including
\begin{description}
  \item [Grid search:] Explore pairwise misfit slices and 3-D misfit surface. Identify local vs global minima.
  \item [Inversion path:] Run Gauss--Newton with different $\lambda$; examine step length and stability.
  \item [Effect of noise:] Increase $\sigma$; observe widening misfit basins and parameter uncertainty.
  \item [Non-uniqueness:] Demonstrate parameter trade-offs (e.g.\ $a^3\, \Delta\rho$) and resolution ellipses.
  \item [Role of Noise:] Look at how noise affects results.
\end{description}

\timeline[slot]{Wrap-up}{10}

Forward models compute fields from distributions of physical properties.  Inverse problems attempt to solve for the distribution of physical properties given a set of observations of a field.

Inversions come with several challenges: non-linearity, non-uniqueness, instability, and underdetermined.  Non-linear functions can sometimes be linearized and if not can be solved using iterative methods.  Just as multiple objects can cast the same shadow, multiple distributions of a property can produce the same geophysical field anomaly.  Damping can help reduce instability and adding geologic constraints can reduce non-uniqueness.

\timeline[slot]{Preview: Gravity Theory}{10}

This week we treated gravity primarily as data to be inverted. Next week we will begin developing gravity theory itself.

We will derive the gravity field from Newton's law of universal gravitation and examine how mass distributions create gravity anomalies. We will see how simple geological bodies such as spheres, cylinders, slabs, and contacts produce characteristic gravity signatures.

Along the way, we will connect the mathematical concepts introduced in Week 1—gradients, curvature, and potential fields—to one of the most important geophysical methods.

By the end of next week, you will be able to predict how changes in density, geometry, and observation distance affect the gravity field and begin interpreting gravity anomalies in geological terms.

% ----------------------------------------------------------
% \section*{References (suggested for students)}
% \begin{thebibliography}{9}
% \bibitem{Blakely} R. Blakely, \textit{Potential Theory in Gravity and Magnetic Applications}, Cambridge University Press, 1995.
% \bibitem{Telford} W. M. Telford, L. P. Geldart, and R. E. Sheriff, \textit{Applied Geophysics}, 2nd ed., Cambridge University Press, 1990.
% \bibitem{Menke} W. Menke, \textit{Geophysical Data Analysis: Discrete Inverse Theory}, 3rd ed., Academic Press, 2012.
% \bibitem{Aster} R. C. Aster, B. Borchers, and C. H. Thurber, \textit{Parameter Estimation and Inverse Problems}, 3rd ed., Elsevier, 2019.
% \bibitem{Tarantola} A. Tarantola, \textit{Inverse Problem Theory and Methods for Model Parameter Estimation}, SIAM, 2005.
% \end{thebibliography}

\workshopend
\notepage
\notepage

\subimport{activities/}{activity_linear_test}
\subimport{activities/}{activity_linearity}
\subimport{activities/}{activity_buried_sphere}
